% =============================================================================
% 文件: sections/abstract.tex   摘要专用页
% 说明: 全文数字由 code/s10_fill_numbers.py 从实算结果自动回填, 不用手抄。
% =============================================================================
\begin{abstract}

中药材热风烘干是多孔介质内耦合传热传质过程，其核心困难在于物性的强非线性、
表面边界值的正确重构以及长达数天的时间尺度。本文把药材简化为长圆柱，
建立了一维轴对称耦合传热传质模型，并用一套有限体积求解器统一求解四个问题。

\textbf{模型方面}，控制方程为径向 Fick 扩散方程
$\partial C/\partial t=(1/r)\partial_r(rD\partial_r C)$ 与含变物性的导热方程
$\rho(C)c_p(C)\partial T/\partial t=(1/r)\partial_r(rk\partial_r T)$，
$D$ 同时依赖含水率与热力学温度（Arrhenius 形式），
表面为第三类对流边界条件（$-k\partial T/\partial r=h(\Tair-\Ts)$、
$-D\partial C/\partial r=h_m(\Cs-\Cair)$），中心对称，初始温度 28 $^\circ$C、
干基含水率 2.55 kg/kg。针对问题 4 的收缩，采用物质坐标 $\xi=r/R(t)$
把变动域映射为固定区间，并证明了该坐标下时间偏导数恰为物质导数、
附加对流项严格为零。数值上采用节点中心有限体积离散、界面扩散系数取调和平均、
表面用"半单元扩散阻力与表面膜阻力串联"重构真实表面值，
时间推进用隐式 Euler 配合固定右端项（旧时间层）的 Picard 迭代。

\textbf{结果方面}，问题 1（1800 s）得到中心温度 33.5754 $^\circ$C、
表面 36.7856 $^\circ$C，中心含水率 2.5500、表面 1.5102 kg/kg，
即 30 min 内热量尚未穿透药材而表面已失水 40.8\%；问题 2（3 h）得到
中心与表面温度 49.8495 与 49.9664 $^\circ$C、含水率 1.7662 与 1.0081 kg/kg；
问题 3 的烘干时间为 \text{57.48}\ h；问题 4 考虑实测收缩后为
\text{51.10}\ h，比不收缩情形缩短约 \text{6.38}\ h，
原因是半径缩短使特征扩散时间中的 $R^2$ 因子减少约 64\%（半径本身减少约 40\%），
且附录 4 的扩散系数随含水率衰减更慢。
四个结果文件均按附件 3 模板给出四位小数。

\textbf{检验方面}，本文建立了五条独立证据链：
（i）与圆柱 Robin 问题的 Bessel 级数解析解对照，空间误差按 $O(N^{-2})$ 收敛
（二阶），并定量证明若把最后一个单元中心当作表面值会产生约
\text{2.28e-03} 的假误差；
（ii）$\Delta t$ 由 10 s 加密到 0.02 s 实测收敛阶 \text{1.00}，
Richardson 外推残余偏差仅 \text{5.40e-05}，
说明隐式 Euler 的一阶误差在 $\Delta t=1\ \mathrm{s}$ 时会让四位小数失真；
（iii）三个算例的水量收支闭合误差均在 $10^{-11}$ 以下（机器精度量级），
含水率非负单调、温度不越界；
（iv）Crank--Nicolson 与隐式 Euler 收敛到同一极限；
（v）合成数据回归（twin experiment）反演传质系数，
真值落在 95\% 置信区间内，残差标准差与注入噪声之比接近 1，
说明模型结构无系统性偏差；并进一步给出可辨识性结论——30 min 窗口的含水率
数据在原理上无法把 $h_m$ 与 $D$ 分开标定（两参数相关系数 \text{+0.827}），
观测窗延长至 24 h 后相关系数降至 \text{-0.096} 才可分别辨识。

\textbf{不确定度方面}，灵敏度分析表明 $h_m$ 与 $D$ 是主控参数，风温的影响很小，
因此提高干燥效率应优先强化对流传质而非升温；
最大的结构不确定源是 $t>14400\ \mathrm{s}$ 的环境平台假设
（附件 1 只覆盖前 4 h），五种口径给出的烘干时间极差为
\text{0.356}\ h，本文以实测末段均值作为基准并给出了区间。
检验还定量揭示了附件 2 的实测收缩轨迹与附录 4 密度经验式之间的不自洽
（干物质质量相对变化 \text{-8.91\%}），据此把问题 4 严格表述为
"给定实测收缩轨迹下的预测"。

\keywords{热风干燥；耦合传热传质；有限体积法；物质坐标与收缩；
半单元表面重构；参数反演与不确定度量化}
\end{abstract}
